\section*{الفصل \ref{ch:g11:drift-to-spreading} --- من الانجراف القاري إلى توسع قاع المحيط}

\begin{solution}{exo:g11:drift-to-spreading:1}
\omterm{def:g11:drift-to-spreading:spreading}{توسع قاع المحيط} هو صنع قاع محيطي جديد عند قمة ظهرة
وسط محيطية، حيث يصعد صخر منصهر ويتصلب، ثم
ابتعاده عن الظهرة على الجانبين. والمعدل الكلي هو السرعة التي
تتباعد بها الصفيحتان؛ ونصف المعدل هو السرعة التي يبتعد بها القاع على
أحد الجانبين عن المحور. وفي التوسع المتماثل، المعدل الكلي $= 2
\times$ نصف المعدل.
\end{solution}

\begin{solution}{exo:g11:drift-to-spreading:2}
فوق \omterm{def:g11:drift-to-spreading:remanence}{درجة حرارة كوري}، أي نحو \qty{580}{\celsius}، لا يستطيع المغنتيت أن يحتفظ
بمغنطة خاصة به؛ ولا تتمغنط حبيباته على امتداد المجال وتحتفظ بذلك الاتجاه
إلا حين يبرد دونها. وإذا سُخّن الصخر من جديد إلى
\qty{700}{\celsius}، فوق \omterm{def:g11:drift-to-spreading:remanence}{درجة حرارة كوري}، فقد سجله القديم؛
وحين يبرد مرة أخرى يسجل مجال ذلك الزمن الجديد.
\end{solution}

\begin{solution}{exo:g11:drift-to-spreading:3}
ينزل الصوت ثم يعود: $d = 1500 \times 5.2/2 = \qty{3900}{m}$.
\end{solution}

\begin{solution}{exo:g11:drift-to-spreading:4}
بدأت الفترة العادية الحالية قبل \num{0.78} مليون سنة: فقد دامت
\num{780000} سنة. ودام حدث هاراميو $1.07 - 0.99 = \num{0.08}$
مليون سنة، أي \num{80000} سنة. وبعدّ الحدود التي يقل عمرها عن خمسة
ملايين سنة (0.78، 0.99، 1.07، 1.77، 1.95، 2.14، 2.15، 2.581، 3.04، 3.11،
3.22، 3.33، 3.58، 4.18، 4.29، 4.48، 4.62، 4.80، 4.89، 4.98): عشرون انقلابًا.
\end{solution}

\begin{solution}{exo:g11:drift-to-spreading:5}
$65\,000/40\,000 \approx 1.6$: تلتف الظهرات حول الأرض نحو مرة
ونصف.
\end{solution}

\begin{solution}{exo:g11:drift-to-spreading:6}
$20/0.78 \approx \qty{25.6}{km}$ في كل مليون سنة، أي نصف معدل
قدره نحو \qty{2.6}{cm} في السنة؛ والمعدل الكلي نحو \qty{5.1}{cm} في السنة.
\end{solution}

\begin{solution}{exo:g11:drift-to-spreading:7}
الغرب: $45/3.58 \approx 12.6$~\unit{km} في كل مليون سنة، أي \qty{1.26}{cm} في السنة؛
الشرق: $44/3.58 \approx 12.3$~\unit{km} في كل مليون سنة، أي \qty{1.23}{cm} في السنة.
المعدل الكلي: $1.26 + 1.23 \approx \qty{2.5}{cm}$ في السنة. ويختلف نصفا المعدل
بنحو \qty{2}{\percent}: فالتوسع متماثل تقريبًا تمامًا.
\end{solution}

\begin{solution}{exo:g11:drift-to-spreading:8}
\qty{1.9}{cm} في السنة تساوي \qty{19}{km} في كل مليون سنة: $990/19 \approx 52$
مليون سنة، أي أقدم قليلًا من عمر أقدم الرواسب البالغ 49 مليون سنة
(بنحو \qty{6}{\percent}).
\end{solution}

\begin{solution}{exo:g11:drift-to-spreading:9}
بمعدل \qty{20}{km} في كل مليون سنة: \qty{15.6}{km}، و\qty{19.8}{km}،
و\qty{21.4}{km}، و\qty{35.4}{km}، و\qty{39.0}{km}. وابتداءً من المحور: عادية من 0
إلى \qty{15.6}{km}، ثم معكوسة حتى \qty{19.8}{km}، ثم عادية حتى \qty{21.4}{km}
(هاراميو)، ثم معكوسة حتى \qty{35.4}{km}، ثم عادية حتى \qty{39.0}{km} (أولدوفاي)،
ثم معكوسة.
\end{solution}

\begin{solution}{exo:g11:drift-to-spreading:10}
$64 \times 10^6/4.6 \times 10^9 \approx \num{0.014}$، أي نحو
\qty{1.4}{\percent} من عمر الأرض. وهو يدعم حجة هس القائلة إن
قاع المحيط حديث لأنه يُصنع ويُدمَّر دائمًا: فلم يبقَ أي
قاع محيطي من معظم تاريخ الأرض.
\end{solution}

\begin{solution}{exo:g11:drift-to-spreading:11}
في كل لحظة يتكوّن القاع عند المحور بقطبية تلك اللحظة،
ثم ينشطر: فيبتعد نصفه على كل جانب بالسرعة نفسها. لذا فإن شريطين تكوّنا
في الزمن نفسه يقعان على البعد نفسه من المحور، واحد على كل
جانب، وبالقطبية نفسها: فكل جانب نسخة من الآخر، منعكسة حول
المحور. وينكسر التماثل إذا ابتعد القاع على أحد الجانبين أسرع منه
على الآخر (توسع غير متماثل)، أو إذا قفز المحور إلى
موضع جديد.
\end{solution}

\begin{solution}{exo:g11:drift-to-spreading:12}
(1) \qty{5}{cm} تساوي $5 \times 10^{-5}$~\unit{km}: فكل سنة يضاف شريط مساحته
$65\,000 \times 5 \times 10^{-5} = \qty{3.25}{km^2}$.
(2) $2.9 \times 10^8/3.25 \approx 9 \times 10^7$ سنة، أي نحو 90 مليون
سنة. (3) وهذا من رتبة عمر أقدم قاع نفسها
(190 مليون سنة) وأقصر منه: فلكي يبقى قاع المحيط بهذه
الحداثة، يجب أن يُدمَّر في مكان ما من القاع قدرُ ما يُصنع عند الظهرات.
\end{solution}

\begin{solution}{exo:g11:drift-to-spreading:13}
دام الحدث \num{0.08} مليون سنة. فبمعدل \qty{12.5}{km} في كل مليون سنة
يكون عرض شريطه $0.08 \times 12.5 = \qty{1}{km}$؛ وبمعدل \qty{75}{km} في كل
مليون سنة، $0.08 \times 75 = \qty{6}{km}$. ومن سطح البحر، على ارتفاع
\qty{3}{km} فوقه، لا يُحدث شريط أضيق بكثير من عمقه إلا مجالًا ضعيفًا،
يتداخل مع مجالَي جاريه العريضين ذوي القطبية
المعاكسة؛ أما شريط عرضه ضعف عمقه فيبرز بوضوح.
\end{solution}

\begin{solution}{exo:g11:drift-to-spreading:14}
الموقع 16: $176/11 \approx \qty{16}{m}$ في كل مليون سنة؛ الموقع 20: $72/67
\approx \qty{1.1}{m}$ في كل مليون سنة. تتراكم الرواسب بمعدلات شديدة
التفاوت من مكان إلى مكان ومن زمن إلى زمن، فلا يعطي سمكها
عمر القاع. وإنما يعطي العمرَ أحافيرُ أقدم
الرواسب، الواقعة مباشرة على البازلت.
\end{solution}

\begin{solution}{exo:g11:drift-to-spreading:15}
(1) الغرب: $58/2.581 \approx \qty{22.5}{km}$ في كل مليون سنة، أي
\qty{2.25}{cm} في السنة؛ الشرق: $70/2.581 \approx \qty{27.1}{km}$ في كل مليون
سنة، أي \qty{2.71}{cm} في السنة؛ والمعدل الكلي $128/2.581 \approx \qty{49.6}{km}$ في كل
مليون سنة، أي نحو \qty{5.0}{cm} في السنة. (2) نحو \qty{225}{km} غربًا
و\qty{271}{km} شرقًا. (3) كان الحدان معًا عند المحور حين وقع
الانقلاب؛ ومنذئذ تباعدا بالمعدل الكلي، أيًّا كان توزيع
الحركة بين الجانبين. فالمسافة بينهما مقسومة على
عمرهما هي المعدل الكلي.
\end{solution}

\begin{solution}{pb:g11:drift-to-spreading:1}
\textbf{1.} الموقع 16 هو الأقرب (\qty{191}{km}، 11 مليون سنة)، والموقع 20
الأبعد (\qty{1270}{km}، 67 مليون سنة): كلما ابتعدنا عن الظهرة،
كان القاع أقدم.

\textbf{2.} تبدأ الرواسب بالترسب على البازلت ما إن يتكوّن
البازلت، فتكون أقدم طبقة في عمر البازلت تقريبًا. ولا يمكن أن تكون إلا
أحدث قليلًا، لا أقدم أبدًا.

\textbf{3.} بالكيلومترات في كل مليون سنة: الموقع 16: $191/11 \approx 17.4$؛
الموقع 15: $380/24 \approx 15.8$؛ الموقع 18: $506/26 \approx 19.5$؛ الموقع 17:
$643/33 \approx 19.5$؛ الموقع 14: $727/40 \approx 18.2$؛ الموقع 19: $990/49
\approx 20.2$؛ الموقع 20: $1270/67 \approx 19.0$.

\textbf{4.} $129.5/7 \approx \num{18.5}$~\unit{km} في كل مليون سنة، أي نحو
\qty{1.85}{cm} في السنة.

\textbf{5.} نعم: نسبتاهما (19.5) تقعان بين نسب المواقع الغربية.
فقد توسع القاع بالمعدل نفسه تقريبًا على الجانبين: والتوسع
متماثل.

\textbf{6.} $t = d/u$: العمر متناسب مع المسافة.

\textbf{7.} $u = 208\,276/10\,952 \approx \num{19.0}$~\unit{km} في كل مليون سنة،
أي \qty{1.90}{cm} في السنة.

\textbf{8.} نحو $2 \times 1.9 = \qty{3.8}{cm}$ في السنة.

\textbf{9.} نعم: \qty{1.9}{cm} في السنة أقل قليلًا من \qty{2}{cm} في
السنة.

\textbf{10.} $1270/19.0 \approx 67$ مليون سنة، وهو عمر أقدم
رواسبه نفسه.

\textbf{11.} $191/9 \approx \num{21.2}$~\unit{km} في كل مليون سنة، أي
\qty{2.1}{cm} في السنة.

\textbf{12.} $1270/71 \approx \num{17.9}$~\unit{km} في كل مليون سنة، أي
\qty{1.8}{cm} في السنة.

\textbf{13.} $180\,329/9854 \approx \num{18.3}$~\unit{km} في كل مليون سنة، أي
\qty{1.83}{cm} في السنة.

\textbf{14.} $(1.90 - 1.83)/1.83 \approx \qty{4}{\percent}$: يتفق المعدلان
في حدود عدم اليقين في الأعمار والمسافات.

\textbf{15.} أحد العمرين مصدره الأحافير في الرواسب، والآخر مصدره
مغناطيسية البازلت والسلم الزمني للانقلابات. والخطأ في إحدى
الطريقتين لا يظهر في الأخرى: فاتفاقهما اختبار قوي
\omterm{def:g11:drift-to-spreading:spreading}{لتوسع قاع المحيط}.

\textbf{16.} $13.8 + 13.8 = \ang{27.6}$.

\textbf{17.} $111.2 \times \cos 30^\circ \approx \qty{96.3}{km}$ لكل درجة؛
$27.6 \times 96.3 \approx \qty{2660}{km}$.

\textbf{18.} $2660/19.0 \approx 140$ مليون سنة.

\textbf{19.} $2660/130 \approx \num{20.5}$~\unit{km} في كل مليون سنة، أي نحو
\qty{2.0}{cm} في السنة، وهو أكثر قليلًا من \qty{1.9}{cm} في السنة. وبما أن آخر
67 مليون سنة جرت بنحو \qty{1.9}{cm} في السنة، فقد كان الانفتاح المبكر
أسرع قليلًا: فالقاع بين الموقع 20 والحافة، $2660 - 1270 =
\qty{1390}{km}$، تكوّن في $130 - 67 = 63$ مليون سنة، أي نحو
\qty{2.2}{cm} في السنة.

\textbf{20.} \textbf{يبلغ عمر أقدم قاع محيطي بجوار إفريقيا عند \ang{30} جنوبًا
نحو 130 مليون سنة} (140 مليون سنة باستقراء المعدل
الحالي): فقد بدأ جنوب الأطلسي ينفتح هناك في العصر الطباشيري.
\end{solution}
